%% bab3.tex
%% BAB 3 - METODOLOGI PENELITIAN
%% Format sesuai Panduan Penyusunan Tesis ITS 2021

\chapter{METODOLOGI PENELITIAN}

Pada bagian ini diuraikan desain, metode, atau pendekatan yang akan digunakan dalam menjawab permasalahan penelitian/studi untuk mencapai tujuan penelitian, serta tahapan penelitian secara rinci, singkat dan jelas.

\section{Desain Penelitian}

Uraikan desain penelitian yang digunakan. Jelaskan apakah penelitian ini bersifat kualitatif, kuantitatif, eksperimental, kajian pustaka, atau kerja pengembangan (\textit{project}).

\section{Alur Penelitian}

Gambar~\ref{fig:alur-penelitian} menunjukkan alur penelitian yang akan dilakukan.

\begin{figure}[htbp]
    \centering
    % Ganti dengan gambar alur penelitian Anda
    \fbox{\parbox{0.8\textwidth}{\centering\vspace{3cm}[Diagram Alur Penelitian]\vspace{3cm}}}
    \caption{Alur Penelitian}
    \label{fig:alur-penelitian}
\end{figure}

\section{Pengumpulan Data}

Jelaskan teknik/metode perolehan data yang digunakan dalam penelitian ini. Uraikan parameter penelitian dan sumber data yang digunakan.

\subsection{Data Primer}

Jelaskan data primer yang digunakan dalam penelitian ini.

\subsection{Data Sekunder}

Jelaskan data sekunder yang digunakan dalam penelitian ini.

\section{Metode yang Digunakan}

Uraikan metode atau pendekatan yang digunakan dalam penelitian ini. Jelaskan bagaimana persoalan yang ada diselesaikan, apakah menggunakan simulasi, model analitis, prototipe, atau analisis data aktual.

\subsection{Metode Pertama}

Jelaskan metode pertama yang digunakan.

\subsection{Metode Kedua}

Jelaskan metode kedua yang digunakan.

\section{Analisis Data}

Jelaskan teknik analisis data yang digunakan dalam penelitian ini.

\section{Evaluasi}

Jelaskan metode evaluasi yang akan digunakan untuk mengukur keberhasilan penelitian.

\section{Jadwal Kegiatan Penelitian}

Tabel~\ref{tab:jadwal} menunjukkan jadwal kegiatan penelitian dalam bentuk \textit{bar-chart}.

\begin{table}[htbp]
    \caption{Jadwal Kegiatan Penelitian}
    \label{tab:jadwal}
    \centering
    \begin{tabular}{|l|c|c|c|c|c|c|}
        \hline
        \textbf{Kegiatan} & \textbf{Bln 1} & \textbf{Bln 2} & \textbf{Bln 3} & \textbf{Bln 4} & \textbf{Bln 5} & \textbf{Bln 6} \\
        \hline
        Studi Literatur & $\bullet$ & $\bullet$ & & & & \\
        \hline
        Pengumpulan Data & & $\bullet$ & $\bullet$ & & & \\
        \hline
        Perancangan Sistem & & & $\bullet$ & $\bullet$ & & \\
        \hline
        Implementasi & & & & $\bullet$ & $\bullet$ & \\
        \hline
        Pengujian & & & & & $\bullet$ & $\bullet$ \\
        \hline
        Penulisan Tesis & $\bullet$ & $\bullet$ & $\bullet$ & $\bullet$ & $\bullet$ & $\bullet$ \\
        \hline
    \end{tabular}
\end{table}
